\documentclass[12pt,a4paper]{article}
\usepackage[utf8]{inputenc}
\usepackage[ngerman]{babel}
\usepackage{amsmath,amssymb}
\usepackage{booktabs}
\usepackage{geometry}
\usepackage{hyperref}
\usepackage{enumitem}
\usepackage{fancyvrb}
\usepackage{mdframed}
\usepackage{parskip}

\geometry{margin=2.5cm}

\title{\textbf{HILOMOT -- Einführung, Motivation und Ansatz}\\\large Eine Lernunterlage für die Hochschule}
\author{}
\date{}

\begin{document}
\maketitle
\tableofcontents
\newpage

%------------------------------------------------------------------
\section{Ausgangsproblem: Warum brauchen wir lokale Modelle?}

In der Praxis begegnen uns ständig \textbf{nichtlineare Systeme} -- ob Verbrennungsmotoren, chemische Reaktoren oder Roboterkinematiken. Ein einfaches lineares Modell $y = a \cdot x + b$ kann solche Zusammenhänge nicht beschreiben.

\textbf{Idee:} Statt ein einziges kompliziertes nichtlineares Modell zu bauen, zerlegen wir das Problem in \textbf{viele kleine Teilprobleme}, die jeweils einfach (linear) lösbar sind.

\begin{mdframed}
\textbf{Analogie:} Stell dir eine Landkarte vor. Statt die gesamte Erdoberfläche mit einer einzigen Formel zu beschreiben, verwenden wir viele kleine, flache Kartenblätter. Jedes Blatt ist für sich genommen einfach (flach/linear), aber zusammen beschreiben sie die gekrümmte Oberfläche beliebig genau.
\end{mdframed}

Das führt uns zum Konzept der \textbf{lokalen Modellnetze:}

\begin{equation}
\hat{y}(x) = \sum_{i=1}^{M} \Phi_i(x) \cdot \hat{y}_i(x)
\end{equation}

\begin{itemize}
  \item $\hat{y}_i(x)$: ein einfaches lineares Modell, das in einer bestimmten Region \glqq zuständig\grqq{} ist
  \item $\Phi_i(x)$: eine Gewichtungsfunktion, die angibt, \textbf{wie stark} Modell~$i$ an der Stelle~$x$ beiträgt
  \item Es gilt immer: $\sum_i \Phi_i(x) = 1$ -- die Gewichte summieren sich zu~1 (Partition der Eins)
\end{itemize}

Die zentrale Frage ist nun: \textbf{Wie teilen wir den Eingangsraum sinnvoll auf?}

%------------------------------------------------------------------
\section{LOLIMOT -- Der bewährte Ansatz}

LOLIMOT (\textit{Local Linear Model Tree}) beantwortet diese Frage mit einer einfachen Strategie:

\textbf{Schneide den Raum immer entlang einer Achse durch.}

Das bedeutet: Jeder Split teilt eine Region in zwei Hälften -- immer parallel zu einer Koordinatenachse. Wie beim Schneiden eines Kuchens, aber nur gerade Schnitte, immer horizontal oder vertikal.

\subsection{Vorgehen (vereinfacht)}
\begin{enumerate}
  \item Starte mit einem einzigen Modell für den ganzen Raum
  \item Finde die Region mit dem größten Fehler
  \item Teile sie in der Mitte -- probiere jede Achse durch, nimm den besten Schnitt
  \item Trainiere die lokalen Modelle neu
  \item Wiederhole, bis der Fehler klein genug ist
\end{enumerate}

\textbf{Stärken:} Einfach, schnell, robust, gut verstanden.

\textbf{Aber:} Was passiert, wenn die \glqq wahre\grqq{} Grenze zwischen zwei Verhaltensbereichen \textbf{schräg} verläuft?

\begin{mdframed}
\textbf{Beispiel:} Ein Motor verhält sich unterschiedlich je nach Kombination von Drehzahl und Last. Die Grenze zwischen \glqq magerem\grqq{} und \glqq fettem\grqq{} Betrieb verläuft diagonal im Drehzahl-Last-Kennfeld. LOLIMOT muss diese diagonale Grenze durch viele kleine Treppenstufen (achsenparallele Schnitte) annähern -- das kostet \textbf{unnötig viele Teilmodelle}.
\end{mdframed}

%------------------------------------------------------------------
\section{HILOMOT -- Die Weiterentwicklung}

\subsection{Kernidee: Schräg schneiden dürfen}

\textbf{HILOMOT} (\textit{Hierarchical Local Model Tree}) löst genau dieses Problem. Statt nur achsenparallel zu schneiden, darf die Trennlinie \textbf{beliebig orientiert} sein:

\begin{center}
\begin{tabular}{lll}
\toprule
 & \textbf{LOLIMOT} & \textbf{HILOMOT} \\
\midrule
\textbf{Schnittrichtung} & Nur entlang der Achsen & Beliebig orientiert (schräg) \\
\textbf{Trennfläche} & $x_j = \text{const}$ & $w^\top x + b = 0$ (Hyperebene) \\
\bottomrule
\end{tabular}
\end{center}

\begin{mdframed}
\textbf{Analogie:} LOLIMOT schneidet den Kuchen nur waagrecht oder senkrecht. HILOMOT darf auch diagonal schneiden -- und trifft damit die natürliche Form des Kuchens oft viel besser.
\end{mdframed}

\subsection{Der Baum als Organisationsstruktur}

HILOMOT organisiert die Aufteilung als \textbf{Binärbaum:}

\begin{Verbatim}[frame=single]
            [Wurzel-Gate]
           /             \\
     [Gate]              Modell 3
     /    \\
Modell 1   Modell 2
\end{Verbatim}

\begin{itemize}
  \item \textbf{Innere Knoten} = Entscheidungen: \glqq Gehört der Punkt eher links oder rechts?\grqq
  \item \textbf{Blätter} = Lokale lineare Modelle
\end{itemize}

An jedem inneren Knoten sitzt eine \textbf{weiche Entscheidungsfunktion} (Sigmoid/Logistik):

\begin{equation}
g(x) = \frac{1}{1 + e^{-(w^\top x + b)}}
\end{equation}

Diese Funktion gibt einen Wert zwischen 0 und 1 aus:
\begin{itemize}
  \item Nahe 1 $\rightarrow$ der Punkt gehört \glqq eher links\grqq
  \item Nahe 0 $\rightarrow$ der Punkt gehört \glqq eher rechts\grqq
  \item Genau 0{,}5 $\rightarrow$ der Punkt liegt auf der Trennfläche
\end{itemize}

\textbf{Wichtig:} Der Übergang ist \textbf{weich/glatt}, nicht hart. Es gibt keinen abrupten Sprung an der Grenze -- das sorgt für ein glattes Gesamtmodell.

\subsection{Wie entsteht das Gesamtgewicht eines Blattes?}

Das Gewicht eines Blattes ergibt sich aus dem \textbf{Produkt aller Entscheidungen auf dem Weg von der Wurzel zum Blatt:}

\begin{mdframed}
Für Modell~1 im Beispielbaum oben:
\begin{itemize}
  \item An der Wurzel: \glqq nach links\grqq{} $\rightarrow$ Gewicht $g_{\text{Wurzel}}(x)$
  \item Am zweiten Knoten: \glqq nach links\grqq{} $\rightarrow$ Gewicht $g_{\text{Knoten}}(x)$
  \item Gesamt: $\Phi_1(x) = g_{\text{Wurzel}}(x) \cdot g_{\text{Knoten}}(x)$
\end{itemize}
\end{mdframed}

Dieses Pfadprodukt garantiert automatisch, dass sich alle Blattgewichte zu~1 summieren.

\subsection{Wie wird HILOMOT trainiert?}

Das Training folgt dem gleichen \textbf{Wachstumsprinzip} wie LOLIMOT:

\begin{enumerate}
  \item \textbf{Starte einfach:} Ein einziges lineares Modell für alle Daten
  \item \textbf{Finde das schwächste Blatt:} Wo ist der Fehler am größten?
  \item \textbf{Spalte es auf:} Ersetze das Blatt durch einen neuen Knoten mit zwei Kindern
  \item \textbf{Optimiere die Trennfläche:} Finde die beste Orientierung $(w, b)$ des Schnitts -- hier liegt der Unterschied zu LOLIMOT:
  \begin{itemize}
    \item LOLIMOT: Probiere $p$~Achsen durch (diskret, einfach)
    \item HILOMOT: Optimiere Richtung und Position kontinuierlich (z.\,B. Gradientenabstieg)
  \end{itemize}
  \item \textbf{Trainiere die lokalen Modelle} (gewichtete kleinste Quadrate)
  \item \textbf{Wiederhole} bis der Fehler akzeptabel ist
\end{enumerate}

%------------------------------------------------------------------
\section{Gegenüberstellung: Wann nehme ich was?}

\begin{center}
\begin{tabular}{lll}
\toprule
\textbf{Aspekt} & \textbf{LOLIMOT} & \textbf{HILOMOT} \\
\midrule
Schnitte & Achsenparallel & Beliebig orientiert \\
Anzahl Modelle & Tendenziell mehr & Oft weniger nötig \\
Training & Schnell, einfach & Aufwendiger \\
Robustheit & Sehr robust & Empfindlicher \\
Interpretierbarkeit & Rechtecke leicht vorstellbar & Schräge Regionen abstrakter \\
Gut geeignet für & Achsennahe Trennungen & Gekoppelte Eingänge \\
\bottomrule
\end{tabular}
\end{center}

\begin{mdframed}
\textbf{Merksatz:} HILOMOT enthält LOLIMOT als Spezialfall. Wenn man die Trennrichtung~$w$ auf Einheitsvektoren~$e_j$ einschränkt, wird jeder HILOMOT-Split zu einem LOLIMOT-Split.
\end{mdframed}

%------------------------------------------------------------------
\section{Zusammenfassung}

\textbf{Beide Verfahren} erzeugen glatte, interpretierbare Modelle aus lokalen linearen Bausteinen. Die Wahl hängt vom konkreten Problem ab:

\begin{itemize}
  \item Wenn die Daten nahelegen, dass die Grenzen zwischen Betriebsbereichen \textbf{schräg} verlaufen, lohnt sich der Mehraufwand von \textbf{HILOMOT}.
  \item Für einen schnellen ersten Ansatz ist \textbf{LOLIMOT} oft die bessere Wahl.
\end{itemize}

\end{document}